\section{Proofs for the exact audit primitives}
\label{sec:static-proofs}

\subsection{Static frontier and post-processing}

\begin{proof}[Proof of Theorem~\ref{thm:exact-existence}]
If a sound $S$ contained $z\in R(\mathcal D)$, an unsafe preimage of $z$ would belong to
$R^{-1}(S)$, a contradiction.  Conversely, a preimage of a point outside $R(\mathcal D)$ cannot be
unsafe, so every subset of $S_{\max}$ is sound.  Since every sound set is contained in
$S_{\max}$, closed $m$-completeness is possible precisely when no $x$ with $h(x)\geq m$ has a code
in $R(\mathcal D)$, which is Eq.~\eqref{eq:closed-completeness}.  Replacing $\geq m$ by $>m$
gives Eq.~\eqref{eq:open-completeness}.
\end{proof}

\begin{proof}[Proof of Proposition~\ref{prop:threshold}]
For every finite $m>\epssep$, Eq.~\eqref{eq:closed-completeness} holds.  For
$0\leq m<\epssep$, the definition of supremum gives an element of $\Gamma_R$ strictly larger than
$m$, so neither completeness condition holds.  At a finite endpoint no element exceeds
$\epssep$, proving open attainability, while closed attainability holds exactly when the supremum
is not itself in $\Gamma_R$.  If the supremum is infinite, $\Gamma_R$ is unbounded.  Taking
infima proves the result.
\end{proof}

The endpoint can fail in both directions one might otherwise conflate.  If a safe boundary state
with $h=0$ aliases an unsafe state, then $\epssep=0$ but closed zero-completeness is impossible.  If
conflicting margins approach one from below without attaining one, closed one-completeness is
attainable even though $\epssep=1$.

\begin{proof}[Proof of Theorem~\ref{thm:data-processing}]
Equality under $R_1$ implies equality after applying $T$, so every conflicting pair for $R_1$ is
also conflicting for $R_2$.  Hence $\Gamma_{R_1}\subseteq\Gamma_{R_2}$; the supremum inequality
and Proposition~\ref{prop:threshold} give the remaining conclusions.
\end{proof}

\subsection{Stochastic separation and post-processing}

\begin{proof}[Proof of Theorem~\ref{thm:stochastic-separation}]
Fix either the closed or open convention and let $H_m$ be its required safe-state set.  Necessity
is pairwise: if $P_x(S)=1$ for $x\in H_m$ and $P_{x'}(S)=0$ for $x'\in\mathcal D$, then
$P_x\perp P_{x'}$.  Conversely, assume every such pair is singular.  For each pair choose
$S_{x,x'}\in\Sigma$ with $P_x(S_{x,x'})=1$ and $P_{x'}(S_{x,x'})=0$.  Countability permits
\[
 S_x=\bigcap_{x'\in\mathcal D}S_{x,x'},\qquad
 S=\bigcup_{x\in H_m}S_x.
\]
Empty intersections are $\Z$ and empty unions are empty.  Countable intersection preserves
$P_x(S_x)=1$, while for every unsafe $x'$, countable subadditivity gives $P_{x'}(S)=0$.
Thus $S$ is the common separator.  The condition is exactly avoidance of the appropriate upper
ray by $\Gamma_P^{\rm st}$, so the supremum and endpoint argument in
Proposition~\ref{prop:threshold} applies verbatim.

If $\Sigma$ separates realized points, then for Dirac laws $\delta_z\perp\delta_{z'}$ exactly when
$z\neq z'$, giving $\Gamma_P^{\rm st}=\Gamma_R$.  Without that qualification two distinct but
measurably indistinguishable codes induce the same law.  Finally, total variation contracts through
every Markov kernel.  Indeed, for a target event $B$, set $f(z)=K(z,B)\in[0,1]$ and use the
layer-cake identity to obtain
\[
 \left|\int f\,dP_x-\int f\,dP_{x'}\right|
 \leq\int_0^1|P_x(f>t)-P_{x'}(f>t)|\,dt
 \leq\|P_x-P_{x'}\|_{\rm TV}.
\]
Taking the supremum over $B$ gives
\[
 \|P_xK-P_{x'}K\|_{\rm TV}\leq\|P_x-P_{x'}\|_{\rm TV}.
\]
Probability laws are mutually singular exactly when their total variation is one.  Hence a
non-singular input pair cannot become singular after post-processing, proving
$\Gamma_P^{\rm st}\subseteq\Gamma_Q^{\rm st}$ and the defect inequality.
\end{proof}

Countability cannot simply be deleted.  On $\Z=[0,1]$, take one required safe state with law
$\delta_t$ for every $t\in[0,1]$ and one unsafe state with the uniform law.  Every safe law is
pairwise singular to the unsafe law, but a separator accepted almost surely by every $\delta_t$
must equal all of $[0,1]$ and therefore has unsafe probability one.

\begin{proof}[Proof of Proposition~\ref{prop:stochastic-tv}]
For any measurable $S$,
\[
 P_x(S^c)+P_{x'}(S)=1-\{P_x(S)-P_{x'}(S)\}.
\]
The Hahn decomposition of the finite signed measure $P_x-P_{x'}$ supplies a measurable set
attaining the bracketed one-sided supremum, which equals total variation for probability laws.
This proves the first identity and its attainability.

For the action statement, let $M_x=P_x\pi$ and $M_y=P_y\pi$ be the induced action laws.  Since
$A_x\cap A_y=\varnothing$,
\[
 M_x(A_x)-M_y(A_x)\geq(1-v_x)-v_y=1-v_x-v_y.
\]
Total-variation contraction through the policy kernel gives
$1-v_x-v_y\leq\|M_x-M_y\|_{\rm TV}\leq\|P_x-P_y\|_{\rm TV}$.
For either relevant pair of code laws $(P,Q)$ and a common prior $\Pi$, the triangle inequality
and Pinsker's inequality give
\[
 \|P-Q\|_{\rm TV}
 \leq\|P-\Pi\|_{\rm TV}+\|Q-\Pi\|_{\rm TV}
 \leq\sqrt{D_{\rm KL}(P\|\Pi)/2}+\sqrt{D_{\rm KL}(Q\|\Pi)/2}.
\]
Substitution into the exact separator identity or the action inequality, followed by clipping a
possibly negative lower bound at zero, proves the two KL statements.
\end{proof}

Every Gaussian law on $\mathbb R^d$ with positive-definite covariance has a strictly positive
Lebesgue density.  Any two such laws are therefore mutually absolutely continuous, and in
particular are not singular.  Degenerate Gaussians can instead be singular on distinct affine
supports, so nondegeneracy is essential.

\subsection{Common-action and randomization details}\label{sec:action-proofs}

\begin{proof}[Proof of Theorem~\ref{thm:det-action}]
The single action $\pi(z)$ satisfies every compatible state's condition exactly when it belongs to
their safe-action intersection.
\end{proof}

\begin{proof}[Proof of Proposition~\ref{prop:action-data-processing}]
For every $z_1\in R_1(\K)$, its fiber $F_{z_1}^{R_1}$ is contained in the
$R_2$ fiber $F_{T(z_1)}^{R_2}$, and intersecting both with the representation-independent set
$\mathcal V_q$ preserves this inclusion.  Hence, for every action $a$ and every nonempty viable
$R_1$ subfiber,
\[
 \inf_{x\in F_{T(z_1)}^{R_2}\cap\mathcal V_q}q(x,a)
 \leq
 \inf_{x\in F_{z_1}^{R_1}\cap\mathcal V_q}q(x,a).
\]
Taking the maximum over actions, negating, applying the positive part, and finally taking the
supremum over $z_1$ gives
$\beta_{\rm CA}^{\mathcal V_q}(R_1)\leq\beta_{\rm CA}^{\mathcal V_q}(R_2)$.
\end{proof}

Let $(\mathcal U,\Sigma_{\mathcal U})$ be a measurable action space and assume every
$\Usafe(x)$ and $\mathcal A_R(z)$ is measurable.  A randomized decision $\mu_z$ is
\emph{strongly safe} when $\mu_z(\mathcal A_R(z))=1$ and pointwise almost surely safe when
$\mu_z(\Usafe(x))=1$ for every fixed $x$ in the fiber.

\begin{theorem}[Qualified randomized action conflict]\label{thm:rand-action}
A strongly safe randomized decision exists if and only if $\mathcal A_R(z)$ is nonempty.  The same
equivalence holds under pointwise almost-sure safety if either the fiber is countable, or the action
space is compact metric, $\mu_z$ is Borel, and every $\Usafe(x)$ is closed.
\end{theorem}

\begin{proof}
A probability measure cannot put mass one on an empty common set, while a common action induces a
safe Dirac measure.  For a countable fiber, the intersection of the probability-one safe-action
events still has probability one.  In the compact-closed case every finite intersection has
probability one and is nonempty.  The closed sets therefore have the finite-intersection property,
so compactness makes their full intersection nonempty.  The converse is again a Dirac measure.
\end{proof}

The qualification is necessary.  Let the action space and an uncountable fiber both be indexed by
$[0,1]$, and set $\Usafe(x)=[0,1]\setminus\{x\}$.  The common set is empty, but the uniform action
distribution is pointwise almost surely safe for every fixed state; it is not strongly safe.

\begin{theorem}[Common-action projection]\label{thm:projection}
Suppose actions lie in $\mathbb R^p$ and every $\Usafe(x)$ in a fixed fiber is closed and convex.
If $\mathcal A_R(z)$ is nonempty, then for every nominal action $u_{\rm nom}(z)$,
\[
 \pi_{\rm LI}(z)=\arg\min_{u\in\mathcal A_R(z)}
 \frac12\|u-u_{\rm nom}(z)\|_2^2
\]
exists uniquely, is uniformly feasible on the fiber, and satisfies
\[
 \langle u_{\rm nom}(z)-\pi_{\rm LI}(z),u-\pi_{\rm LI}(z)\rangle\leq0
 \quad\forall u\in\mathcal A_R(z).
\]
\end{theorem}

\begin{proof}
The common set is a closed convex intersection.  A minimizing sequence has bounded objective and is
therefore bounded; in finite dimensions it has a convergent subsequence whose limit is feasible and
minimizing.  If two distinct minimizers $u,v$ existed, their feasible midpoint would satisfy
\[
 \left\|\frac{u+v}{2}-u_{\rm nom}\right\|^2
 =\frac12\|u-u_{\rm nom}\|^2+
   \frac12\|v-u_{\rm nom}\|^2-\frac14\|u-v\|^2,
\]
a strict improvement.  The variational inequality follows from the right derivative along the
feasible segment from the minimizer to any $u$ in the common set.
\end{proof}

This result does not ensure that $z\mapsto\pi_{\rm LI}(z)$ is measurable or continuous, nor that an
infinite constraint intersection has a finite verified reduction.  Barrier hypotheses needed for
forward invariance remain separate.

\begin{proof}[Proof of Theorem~\ref{thm:helly-action}]
For $r\geq0$ and $x\in F_z^R$, define the relative expansion
\[
 B_x(r)=\{u\in\mathcal U:d_2(u,\Usafe(x))\leq r\}.
\]
It is compact and convex, and for every nonempty $G$,
\[
 \beta_{\rm geo}(G)=\min\{r\geq0:\bigcap_{x\in G}B_x(r)\neq\varnothing\}.
\]
The minimum exists because the supremum of the continuous distance functions is lower
semicontinuous on compact $\mathcal U$.

Let $b$ be the supremum of $\beta_{\rm geo}(G)$ over subfibers of size at most $p+1$.
Monotonicity under set inclusion gives $b\leq\beta_{\rm geo}(F_z^R)$.  If $r>b$, every at-most-
$(p+1)$ subfamily of $\{B_x(r)\}_{x\in F_z^R}$ intersects.  Helly's theorem makes every finite
subfamily intersect; since every member is compact, the finite-intersection property makes the
whole family intersect.  Thus $\beta_{\rm geo}(F_z^R)\leq r$.  Letting $r\downarrow b$ proves the
displayed equality.  At $r=0$, the same argument says that the full common-action set is nonempty
exactly when every at-most-$(p+1)$ intersection is nonempty.  Attainment and closedness make
$\beta_{\rm geo}=0$ equivalent to a nonempty intersection.  A finite fiber has only finitely many
eligible subfibers, so its supremum is a maximum.

Sharpness follows with $B=\{u:\|u\|_2\leq\sqrt p\}$,
\[
 A_i=B\cap\{u:u_i\geq1\}\quad(1\leq i\leq p),\qquad
 A_{p+1}=B\cap\{u:\textstyle\sum_i u_i\leq p-1\}.
\]
The full intersection is empty.  Omitting $A_{p+1}$ leaves $(1,\ldots,1)$ feasible; omitting
$A_j$ leaves the vector with coordinate $j$ zero and every other coordinate one.  Hence every
$p$-subfamily intersects and $p+1$ cannot be improved.
\end{proof}

When $\mathcal U$ has affine dimension $q<p$, an affine isometry identifies its affine hull with
$\mathbb R^q$; the same proof then replaces $p+1$ by $q+1$.

The assumptions delimit the result.  In one dimension the three closed nonconvex sets
$\{0,1\},\{1,2\},\{0,2\}$ intersect pairwise but not jointly.  Without compactness, an empty
intersection need not have positive geometric slack: the ray $\{(t,0):t\geq0\}$ and the closed
convex epigraph $\{(t,s):t>0,\ s\geq1/t\}$ are disjoint at distance zero.  For an infinite
noncompact family, $[n,+\infty)$, $n\geq1$, has every finite intersection nonempty but empty total
intersection.  These examples rule out the corresponding stronger variants.

\subsection{Finite-action slack and dual witness}

\begin{proof}[Proof of Theorem~\ref{thm:finite-action}]
A Dirac mass on a deterministic minimizer gives
$\beta_{\rm exp}\leq\beta_{\rm det}$.  For any $p\in\Delta_k$, choose $a_p$ with
$p_{a_p}\geq1/k$.  Nonnegativity gives
\[
 \sup_x\sum_a p_a\ell(x,a)
 \geq p_{a_p}\sup_x\ell(x,a_p)
 \geq\frac1k\min_a\sup_x\ell(x,a).
\]
Taking the infimum over $p$ proves the sandwich.  The action-wise robust objective is linear in
$p$, so its minimum is attained at an action minimizing the robust loss.  For any nonempty support,
the maximum supported robust loss is at least the minimum across all actions, and a Dirac mass on a
minimizer attains equality; hence $\beta_{\rm ar}=\beta_{\rm ws}=\beta_{\rm det}$.

If the common safe-action set is empty and loss vanishes exactly on safe actions, every action has
strictly positive loss at some compatible state.  Finiteness of the action set makes
$\beta_{\rm det}>0$, and the sandwich propagates positivity.  Hard failure loss makes every
action-wise supremum one.
\end{proof}

When $F=\{x_1,\ldots,x_n\}$ and $L_{ia}=\ell(x_i,a)$, finite linear-program duality gives
\[
 \beta_{\rm exp}(z)=\max_{q\in\Delta_n}\min_a\sum_{i=1}^n q_iL_{ia}.
\]
The primal minimizes $t$ subject to $Lp\leq t\mathbf1$ and $p\in\Delta_k$; its dual maximizes $r$
subject to $L^Tq\geq r\mathbf1$ and $q\in\Delta_n$.  The maximizing $q$ is a finite adversarial
mixture over compatible states, not a population certificate.

\subsection{Smooth localized collision}

\begin{proof}[Proof of Proposition~\ref{prop:average-separation}]
Let $\psi(t)=\exp(1-(1-t^2)^{-1})$ for $|t|<1$ and zero otherwise, and define
\[
 R_\delta(x)=x-a\psi((x-a)/\delta)+a\psi((x+a)/\delta).
\]
The two bump supports are disjoint and interior to $[-1,1]$, while
$R_\delta(a)=R_\delta(-a)=0$.  Thus a safe point of margin $a$ collides with an unsafe point.
With decoder and latent predictor equal to the identity, reconstruction error is supported on
intervals of total length $4\delta$ and has magnitude at most $a$; the uniform density gives the
bound $2a^2\delta$.  Identity dynamics makes latent prediction exact, and decoded one-step
prediction has the same error as reconstruction.  Finally, $a$ admits only action $+1$ and $-a$
only action $-1$ under the condition $ux\geq0$.
\end{proof}
